\documentclass[11pt]{article}
\usepackage[T1]{fontenc}
\usepackage{textcomp}
\usepackage[margin=0.75in]{geometry}
\usepackage{amsmath,amssymb,amsthm}
\usepackage{array,booktabs}
\usepackage{hyperref}
\setlength{\parindent}{0pt}
\newtheorem{theorem}{Theorem}
\theoremstyle{definition}
\newtheorem{definition}{Definition}
\title{Flow Proof Artifact: Ideal-add-closed}
\author{Generated by \texttt{flow doc proof}}
\date{Flow Proof Book}
\begin{document}
\maketitle
Ideal addition closure follows from ideal axioms.\\[0.5em]
\textbf{Source.} dummit-foote — *Abstract Algebra*, §7.3\\[0.5em]
\bigskip
\noindent{\large \textbf{Derived fact 1.} \textit{If a in I and b in I then a + b in I when I is an ideal} \hfill Ideal \cdot addition \cdot ideal sums stay in the ideal}\par
\smallskip\noindent\textit{Coordinate: Ideal \cdot addition \cdot ideal sums stay in the ideal}\par
\smallskip\noindent\textit{Source: dummit-foote}\par
\smallskip\noindent\textit{Built on: closed under addition, for addition on Ideal, zero is the left identity, for addition on Ring}\par
\medskip
\noindent\textbf{Goal.} If a in I and b in I then a + b in I when I is an ideal\par
\noindent$\displaystyle \forall I \in Ideal \forall a \in Ring \forall b \in Ring\quad a + b in I$\par
\medskip
\noindent
\begin{tabular}{@{} >{\raggedright\arraybackslash}p{0.06\textwidth} >{\raggedright\arraybackslash}p{0.47\textwidth} >{\raggedleft\arraybackslash}p{0.41\textwidth} @{}}
\textbf{\#} & \textbf{Proof} & \textbf{Mathematics} \\
\midrule
\textbf{1.} & We prove that ideal sums stay in the ideal for addition on Ideal. & \\[0.45em]
\textbf{2.} & We split into exhaustive cases — the claim must hold in each one. & \\[0.45em]
\textbf{3.} & Case 1 (see step 2): suppose a in I. & \\[0.45em]
\textbf{4.} & Case 2 (see step 2): suppose b in I. & \\[0.45em]
\textbf{5.} & We invoke the definitional clause governing addition on Ideal: closed under addition, for addition on Ideal (instantiated for I, a, b). & \\[0.45em]
\textbf{6.} & We invoke the definitional clause governing addition on Ring: zero is the left identity, for addition on Ring (instantiated for a). & \\[0.45em]
\textbf{7.} & From step 4, step 5, and step 6, this implies a plus b in I. Together with the other cases (step 3 and step 4), the goal is discharged. Hence proven. & $\displaystyle a + b in I$ \\[0.45em]
\bottomrule
\end{tabular}
\label{thm:Ideal-addition-ideal_sums_stay_in_the_ideal}
\par\medskip\hrule
\end{document}
